\documentclass[10pt,a4paper]{amsart}
\usepackage[margin=19mm]{geometry}
\usepackage{amsmath,amssymb,mathtools}
\usepackage[scr=rsfs,cal=euler]{mathalfa}
\usepackage{xfrac}
\usepackage[unicode,hidelinks,bookmarks=false]{hyperref}
\newcommand{\ie}{i.e.\ }
\newcommand{\cf}{cf.\ }
\newcommand{\resp}{respectively\ }
\newcommand{\Subsection}{\subsection}
\providecommand{\nobreakdash}{\nobreak\hbox{-}}
\setlength{\emergencystretch}{2em}
\multlinegap0pt
\usepackage{bm}
\usepackage{bbm}
\usepackage{dsfont}
\usepackage{xspace}
\usepackage{cancel}
\usepackage{mathtools}
\usepackage{amsthm}
\makeatletter
\@ifundefined{theo}{\newtheorem{theo}{Theo}}{}
\@ifundefined{defi}{\newtheorem{defi}[theo]{Definition}}{}
\@ifundefined{lemm}{\newtheorem{lemm}[theo]{Lemma}}{}
\@ifundefined{prop}{\newtheorem{prop}[theo]{Proposition}}{}
\makeatother
\newcommand{\I}{\mathbb{I}}
\newcommand{\N}{\mathbb{N}}
\renewcommand{\geq}{\geqslant}
\renewcommand{\leq}{\leqslant}
\title[Conjecture: the set of prime numbers is supernatural]{Conjecture: the set of prime numbers is supernatural}
\author{Arnaud Mayeux}
\date{Source: arXiv:2412.12041v2 (CC BY 4.0)}
\begin{document}
\maketitle

\noindent\emph{Source and licence.} Conjecture: the set of prime numbers is supernatural. Version arXiv:2412.12041v2 (CC BY 4.0), distributed under the Creative Commons Attribution 4.0 International licence, as recorded in the arXiv licence metadata for this exact version and shown on the abstract page https://arxiv.org/abs/2412.12041v2. Full rights notice: licenses/arxiv-2412.12041-CC-BY-4.0.txt.\par\smallskip

\noindent\textit{Editorial scope.} Counted unit: the Prop whose source label is \texttt{increa} in the article (source0.tex lines 255--257, source label \texttt{increa}), delivered together with the article's own complete original proof of it (source0.tex lines 259--292, one \texttt{proof} environment of 34 lines). No mathematics is written, repaired or re-derived: statement and proof are the article's own text line for line. Required context retained uncounted: none. Every definition, display and label of those lines is retained unchanged. Omitted from this package and not redistributed: 1--254, 258--258, 293--410. Nothing inside a retained excerpt is omitted or altered: every retained body is delivered exactly as the article's lines are written, so no omission receipt of any kind accompanies this package. Citations: the retained excerpts carry no citation command at all, so no bibliography entry of the article is transcribed and no reference is redistributed. Reference adaptation: none was needed and none was made; every reference inside the delivered unit resolves inside it, so no unresolved reference or citation is produced. Printed slips kept literal: no printed word, symbol, hypothesis or number of the counted statement or of its proof was altered. Layout and class: the article's own class is \texttt{alggeom} with packages including mathrsfs, amscd, amsmath, amssymb, latexsym, comment, wasysym, tikz, tikz-cd, appendix, nicefrac, fix-cm, marvosym, xy; the wrapper is the lane's pinned \texttt{amsart} editorial wrapper at 10pt on A4, which restates the theorem-like environments the retained text needs and adds the article's own macro definitions that the retained text needs (\texttt{\textbackslash I}, \texttt{\textbackslash N}, \texttt{\textbackslash geq}, \texttt{\textbackslash leq}), with extra wrapper packages (bm, bbm, dsfont, xspace, cancel, mathtools). The delivered numbering is the unit's own. Assets: no retained span contains a figure environment, a graphics-inclusion command, a PDF-page inclusion, a table or a TikZ picture; no image file is copied into this package, the package declares no asset dependency and no image file is redistributed with it. Licence: version arXiv:2412.12041v2 of this article is distributed under the Creative Commons Attribution 4.0 International licence, as recorded in the arXiv licence metadata for this exact version and shown on the abstract page https://arxiv.org/abs/2412.12041v2. A full rights notice is delivered at licenses/arxiv-2412.12041-CC-BY-4.0.txt. Characters: every retained excerpt is pure ASCII, so the delivered engine, XeLaTeX with its default Latin Modern text font, reports no missing character.
\begin{prop} \label{increa} A natural function $f \in \mathscr{F} _{Natural}$ is constant or strictly increasing.

\end{prop}

\begin{proof} By definition a function $f $ is in $\mathscr{F} _{Natural}$ is and only if there exists a finite word $\sigma$ in letters $A_+ , A_{\times} , A_{\wedge}$ such that $f \in \sigma (\mathscr{F} _s )$.  We need the following definition.

\begin{defi} The length of a natural function $f \in \mathscr{F}_{Natural}$ is the number \[ \min \{n \in \N \mid  f \in \sigma (\mathscr{F}_s) \text{ with $\sigma$ a word of length $n$} \}, \] it is well defined.

\end{defi} Let us prove our result by induction on length. Consider the following assertion.
\begin{center}

 \textit{$P_n :$ Proposition \ref{increa} is true for every $f \in \mathscr{F}_{Natural}$ of length $\leq n$.}
\end{center}
 
 Let us prove by induction that $P_n$ is true for every $n \in \N$. We need the following Lemmas.
 
  \begin{lemm} \label{lem0} Let $h,g \in \mathscr{F}_{Natural}$ be strictly increasing, then $h+g , h\times g , h \wedge g  $ are strictly increasing.
\end{lemm} 
\begin{proof}
Let $i>j$ be integers in $\I$, then $h(i)>h(j)$ and $g(i)> g (j)$. So $h(i)+g(i) > h(j)+g(j)$, consequently $h+g$ is strictly increasing. Since $ g(i) > g(j) \geq 1 $ and $h(i) > h(j) \geq 1$, we have $h(i)\times g(i) > h(j)\times g(j)$ and consequently $h\times g $ is strictly increasing. We also have $h(i)^{g(i)} > h(j) ^{g(j)}$ and consequently $h \wedge g$ is strictly increasing.
\end{proof} 
 
  \begin{lemm} \label{lem1} Let $h \in \mathscr{F}_{Natural}$ be strictly increasing and $g  \in \mathscr{F} _{Natural}$ be constant, then $h+g , h\times g , h \wedge g  $ are strictly increasing.
\end{lemm} 
\begin{proof}
Let $i>j$ be integers in $\I$, then $h(i)>h(j)$ and $g(i)= g (j)$. So $h(i)+g(i) > h(j)+g(j)$, consequently $h+g$ is strictly increasing. Since $g(i) = g(j) \geq 1$ and $h(i) > h(j)$, we have $h(i)\times g(i) > h(j)\times g(j)$ and consequently $h\times g $ is strictly increasing. We have $h(i)^{g(i)} > h(j) ^{g(j)}$ and consequently $h \wedge g$ is strictly increasing.
\end{proof} 

\begin{lemm} \label{lem2} Let $h \in \mathscr{F}_{Natural}$ be strictly increasing and $g  \in \mathscr{F} _{Natural}$ be constant, then $g \wedge h $ is strictly increasing except if $g=1$. If $g=1$, $g \wedge h $ is constant.
\end{lemm}

\begin{proof} Let $i>j$ be integers in $\I$. If $g=1$ then $g(i) ^{ h(i)}=g(j) ^{ h(j)}=1$ and $g \wedge h$ is constant. In the other case $g(i)=g(j) >1$ and $h(i)>h(j)$, consequently $g(i) ^ {h(i)} > g (j) ^{ h(j)}$. This finishes the proof.
\end{proof}
 
 Now let us start our induction. If $n=0$, then $f$ is constant or the identity map and thus satisfies $P_0$. Now assume that $P_n$ is true and let us prove that $P_{n+1}$ is true. Let $f $ of length $n+1$.  By definition, there are $h,g $ of length $\leq n$ such that $f=h+g$ or $f=h \times g $ or $f = g \wedge h$ or $f= h \wedge g$. By induction hypothesis $h$ and $g$ are constant or strictly increasing.  Assume first both are constant. Then $f$ is obviously constant. If both are strictly increasing, then $f$ is strictly increasing by Lemma \ref{lem0}. If $h$ is strictly increasing and $g$ is constant, then by Lemma \ref{lem1}, $h+g$ , $h \times g$ and $h \wedge g$ are strictly increasing, and by Lemma \ref{lem2}, $g \wedge h$ is strictly increasing or constant. If $h$ is constant and $g$ is strictly increasing, we deduce, in the same way, that $h+g$ , $h \times g$, $h \wedge g$ and $g \wedge h$ are strictly increasing or constant. Consequently, $f$ is constant or strictly increasing. This finishes the induction and the proof of Proposition.


\end{proof}

\end{document}
